\section{Distributed Massive MIMO System Model}

\subsection{Downlink System Model}

We consider the downlink of a distributed (cell-free) massive \gls{mimo}
network in which $L$ \glspl{ap}, each equipped with $M$ antennas, jointly serve
$K$ single-antenna users over the same time-frequency resources. The APs are
connected through fronthaul links to a central processing unit (CPU) that
encodes the data and computes the precoders, so that the $LM$ antennas act as a
single distributed array. Transmission uses OFDM with $Q$ data subcarriers,
indexed by $q=1,\dots,Q$; the model below is stated per subcarrier and the same
processing is applied on each of them.

Let $\vect{h}_{kl}[q]\inC{M}$ denote the channel from AP $l$ to user $k$ on
subcarrier $q$. Stacking the contributions of all APs gives the collective
channel
\begin{equation}
\label{eq:collective-channel}
\vect{h}_{k}[q] =
\begin{bmatrix} \vect{h}_{k1}[q] \\ \vdots \\ \vect{h}_{kL}[q] \end{bmatrix}
\inC{LM}.
\end{equation}
The channel realizations are drawn from the 3GPP TR 38.901 urban-microcell
(UMi) geometry-based stochastic channel model \cite{tr38901}, using the
implementation provided by the Sionna link-level simulator
\cite{hoydis2022sionna}. Each AP is modelled as a 38.901 base station carrying
an $M$-element half-wavelength \gls{ula} of omnidirectional elements and each
user as a single-antenna terminal. Every AP-user link is assigned a line-of-sight
or non-line-of-sight condition according to the distance-dependent LOS
probability of the UMi scenario, and its distance-dependent path loss, log-normal
shadow fading, and small-scale multipath then determine $\vect{h}_{kl}[q]$.

Because the model is specified directly in terms of channel realizations rather
than through a closed-form spatial correlation matrix, the large-scale fading
coefficient $\beta_{kl}$ between AP $l$ and user $k$, which drives the downlink
power control of \Cref{ssec:designs}, is obtained from the realizations as
the average per-antenna channel gain,
\begin{equation}
\label{eq:beta-estimate}
\beta_{kl} = \frac{1}{M}\,\expt{\norm{\vect{h}_{kl}[q]}^{2}},
\end{equation}
which is estimated by averaging $\abs{[\vect{h}_{kl}[q]]_{m}}^{2}$ over the $M$
antennas of AP $l$ and the $Q$ subcarriers. It captures the combined effect of
path loss and shadow fading (and, on a LOS link, of the specular component) and
so summarizes the slow, position-dependent part of the channel that the power
allocation acts on.

The CPU maps the data symbol $s_{k}[q]\in\mathbb{C}$ of user $k$, with unit
power $\expt{\abs{s_{k}[q]}^2}=1$ and mutually independent across users, onto the
$M$ antennas of each AP. AP $l$ transmits
\begin{equation}
\label{eq:ap-transmit}
\vect{x}_{l}[q] = \sum_{k=1}^{K} \vect{w}_{kl}[q]\, s_{k}[q]
= \mat{W}_{l}[q]\,\vect{s}[q],
\end{equation}
where $\vect{w}_{kl}[q]\inC{M}$ is the precoding vector that AP $l$ applies to
user $k$, $\vect{s}[q]=[\,s_{1}[q]\ \cdots\ s_{K}[q]\,]^{\intercal}\inC{K}$
collects the data symbols, and
$\mat{W}_{l}[q]=[\,\vect{w}_{1l}[q]\ \cdots\ \vect{w}_{Kl}[q]\,]\inC{M\times K}$
is the precoding matrix of AP $l$. Collecting the per-AP precoders of a given
user into
$\vect{w}_{k}[q]=[\,\vect{w}_{k1}[q]^{\intercal}\ \cdots\ \vect{w}_{kL}[q]^{\intercal}\,]^{\intercal}\inC{LM}$
yields the network-wide precoding matrix
\begin{equation}
\label{eq:network-precoder}
\mat{W}[q] = \big[\,\vect{w}_{1}[q]\ \cdots\ \vect{w}_{K}[q]\,\big]
\inC{LM\times K},
\end{equation}
whose $l$-th block of $M$ rows equals $\mat{W}_{l}[q]$, so that the stacked
transmit vector is
$\vect{x}[q]=[\,\vect{x}_{1}[q]^{\intercal}\ \cdots\ \vect{x}_{L}[q]^{\intercal}\,]^{\intercal}=\mat{W}[q]\,\vect{s}[q]\inC{LM}$.

The signal received by user $k$ on subcarrier $q$ is
$y_{k}[q]=\vect{h}_{k}[q]^{H}\vect{x}[q]+n_{k}[q]$, which, using
\eqref{eq:network-precoder}, separates into the desired term, the multi-user
interference, and the noise,
\begin{equation}
\label{eq:downlink-received}
y_{k}[q] =
\underbrace{\vect{h}_{k}[q]^{H}\vect{w}_{k}[q]\, s_{k}[q]}_{\text{desired signal}}
+ \underbrace{\sum_{\substack{i=1 \\ i\neq k}}^{K}
\vect{h}_{k}[q]^{H}\vect{w}_{i}[q]\, s_{i}[q]}_{\text{multi-user interference}}
+ \underbrace{\vphantom{\sum_{i=1}^{K}} n_{k}[q]}_{\text{noise}},
\end{equation}
where $n_{k}[q]\sim\cn{0}{\sigma^{2}}$ is the receiver noise with power
$\sigma^{2}$, independent across users and subcarriers and independent of the
channels and data. The scalar $\vect{h}_{k}[q]^{H}\vect{w}_{k}[q]$ is the
effective channel seen by user $k$: the distributed precoder turns the
$LM$-antenna channel into an equivalent single-antenna link.

Treating the multi-user interference in \eqref{eq:downlink-received} as noise,
the effective downlink \gls{sinr} of user $k$ on subcarrier $q$ is
\begin{equation}
\label{eq:downlink-sinr}
\mathrm{SINR}_{k}^{\mathrm{dl}}[q] =
\frac{\abs{\vect{h}_{k}[q]^{H}\vect{w}_{k}[q]}^{2}}
{\displaystyle\sum_{\substack{i=1 \\ i\neq k}}^{K}
\abs{\vect{h}_{k}[q]^{H}\vect{w}_{i}[q]}^{2} + \sigma^{2}} .
\end{equation}
With the effective channel known at the receiver, an achievable rate for user
$k$ on subcarrier $q$ is $\log_{2}\!\big(1+\mathrm{SINR}_{k}^{\mathrm{dl}}[q]\big)$,
and the delivered downlink sum rate $R_{\mathrm{DL}}$ follows by inserting these
per-user, per-subcarrier SINRs into the ergodic rate expression of
\Cref{sec:rate} (\eqref{eq:rate}), which in turn drives the encoder and MIMO
processing consumption of the power model.


\subsection{Precoder and Power Control Designs}
\label{ssec:designs}

Unlike co-located massive MIMO, where a single sum-power budget applies to the
whole array, each AP is a separate transmitter and is subject to its own power
constraint. Since the data symbols are independent and unit-power, the average
power radiated by AP $l$ across the OFDM block is
\begin{align}
\label{eq:power-constraint}
\sum_{q=1}^{Q}\expt{\norm{\vect{x}_{l}[q]}^{2}}
&= \sum_{q=1}^{Q}\sum_{k=1}^{K}\expt{\norm{\vect{w}_{kl}[q]}^{2}}
\leq \rho_{\max}, \\
  l&=1,\dots,L,
\end{align}
where $\rho_{\max}$ is the maximum downlink transmit power per AP, assumed equal
for all APs, and the expectation is taken over the data and channel
realizations.

It is convenient to split each precoder into a transmit direction and a power
control coefficient,
\begin{equation}
\label{eq:precoder-decomposition}
\vect{w}_{k}[q] = \sqrt{\rho_{k}}\,
\frac{\bar{\vect{w}}_{k}[q]}{\sqrt{\expt{\norm{\bar{\vect{w}}_{k}[q]}^{2}}}},
\end{equation}
where $\bar{\vect{w}}_{k}[q]\inC{LM}$ is an arbitrarily scaled vector that fixes
the spatial direction of the transmission to user $k$ and $\rho_{k}\geq 0$ is the
downlink power allocated to that user, so that
$\expt{\norm{\vect{w}_{k}[q]}^{2}}=\rho_{k}$. This separation isolates two design
tasks, the choice of the precoding directions $\{\bar{\vect{w}}_{k}[q]\}$ and the
power allocation that ties them to the per-AP budget in
\eqref{eq:power-constraint}. How both are computed depends on the level of
cooperation between the APs, so the two operation modes are treated in turn:
centralized operation (\Cref{sssec:centralized}), where the CPU designs all
precoders jointly and assigns one power coefficient per user, and local
operation (\Cref{sssec:local}), where each AP designs its own precoder and sets
its powers from local information alone. In both cases the allocation depends
mainly on the large-scale fading of \eqref{eq:beta-estimate} and can be held
fixed over many coherence blocks, whereas the directions are recomputed every
block.



\subsubsection{Centralized Operation}
\label{sssec:centralized}

In centralized operation the APs forward their received uplink pilots to the CPU,
which estimates the collective channels $\{\vect{h}_{k}[q]\}$ and designs all
precoders jointly \cite{demir2021foundations}. Exploiting TDD reciprocity, these
uplink estimates serve as the downlink \gls{csi}, so no downlink pilots are
needed. Because the CPU sees the whole $LM$-antenna array, it can make the
serving APs cancel one another's interference, using up to $LM$ spatial degrees
of freedom per user. Collecting the user channels into
$\mat{H}[q]=[\,\vect{h}_{1}[q]\ \cdots\ \vect{h}_{K}[q]\,]\inC{LM\times K}$, the
\gls{zf} precoder picks the directions that null all inter-user interference,
\begin{equation}
\label{eq:zf-precoder}
\bar{\mat{W}}[q] = \mat{H}[q]\big(\mat{H}[q]^{H}\mat{H}[q]\big)^{-1}
\inC{LM\times K},
\end{equation}
whose column $\bar{\vect{w}}_{k}[q]$ is orthogonal to every other user's channel
when the \gls{csi} is accurate; this requires $K\leq LM$ so that the Gram matrix
$\mat{H}[q]^{H}\mat{H}[q]$ is invertible. Zero-forcing disregards the noise, so
at low SNR it is improved by adding a diagonal loading $\lambda$ to the Gram
matrix, giving the regularized family
\begin{equation}
\label{eq:centralized-rzf}
\bar{\mat{W}}[q] = \mat{H}[q]\big(\mat{H}[q]^{H}\mat{H}[q]+\lambda\mat{I}_{K}\big)^{-1},
\end{equation}
which trades a little residual interference for robustness. With the loading set
to the noise power, $\lambda=\sigma^{2}$, this is the downlink \gls{mmse}
precoder, the dual of \gls{mmse} combining; under the perfect-\gls{csi}
assumption used here the \gls{mmse} precoder coincides with \gls{rzf} at that
loading, and the two differ only once a channel-estimation-error covariance is
modelled. Zero-forcing, the $\lambda=0$ limit, is the scheme adopted in the power
model of \Cref{sec:pwr_model}.

Because a centralized precoder is designed jointly across the APs, a single power
coefficient $\rho_{k}\geq 0$ is attached to each user and is common to all APs;
the direction $\bar{\vect{w}}_{k}[q]$ already fixes how that power is split over
the $LM$ antennas, and the direction is scaled to radiate $\rho_{k}$ as in
\eqref{eq:precoder-decomposition}. Two nominal allocations set $\{\rho_{k}\}$,
both spending the total network budget $P_{\mathrm{tot}}=L\rho_{\max}$. Equal
power gives every user the same share, $\rho_{k}=P_{\mathrm{tot}}/K$, whereas
fractional power control weights users by their aggregate large-scale gain
$\beta_{k}=\sum_{l=1}^{L}\beta_{kl}$,
\begin{equation}
\label{eq:centralized-fractional}
\rho_{k} = P_{\mathrm{tot}}\,
\frac{\beta_{k}^{\,v}}{\sum_{i=1}^{K}\beta_{i}^{\,v}},
\end{equation}
with the exponent $v$ trading throughput ($v>0$, favouring strong
users) against fairness ($v<0$, protecting weak users). These nominal powers
make the per-AP loads average to $\rho_{\max}$ but do not by themselves respect
each per-AP constraint in \eqref{eq:power-constraint}. Since the joint directions
cannot be rescaled per AP without breaking the interference nulling, a single
common factor
\begin{equation}
\label{eq:power-normalization}
\gamma = \sqrt{\frac{\rho_{\max}}{\max_{l}\tilde{p}_{l}}},
\qquad
\tilde{p}_{l} = \sum_{q=1}^{Q}\sum_{k=1}^{K}
\expt{\norm{\vect{w}_{kl}[q]}^{2}},
\end{equation}
rescales all precoders, where $\tilde{p}_{l}$ is the power AP $l$ would radiate
under the nominal allocation. The busiest AP then transmits at exactly
$\rho_{\max}$ and the others stay below their budget, while the common scalar
preserves the relative amplitudes and phases across APs and hence the
interference nulling of \eqref{eq:zf-precoder}. Because $\sum_{k}\rho_{k}=L\rho_{\max}$,
the per-AP powers average to $\rho_{\max}$, so $\gamma\leq 1$ and the constraint
is always feasible.

\subsubsection{Local Operation}
\label{sssec:local}

In local (distributed) operation the CPU only encodes the data, and each AP forms
its own precoding directions from its locally estimated channels
$\hat{\vect{h}}_{kl}[q]$, without sharing instantaneous \gls{csi}. An AP can then
suppress only the interference it itself generates, using the $M$ spatial degrees
of freedom of its own array, so the coherent cross-AP cancellation of the
centralized case is lost. The simplest choice is maximum-ratio transmission
(\gls{mr}), which points each beam along the local channel,
\begin{equation}
\label{eq:mr-precoder}
\bar{\vect{w}}_{kl}[q] = \vect{h}_{kl}[q],
\end{equation}
maximizing the signal power delivered to the desired user but ignoring the
interference it causes. Interference is instead suppressed locally by the
counterpart of zero-forcing applied to each AP's own users. Writing
$\mat{E}_{l}[q]=[\,\vect{h}_{1l}[q]\ \cdots\ \vect{h}_{Kl}[q]\,]\inC{M\times K}$
for AP $l$'s local channel matrix, the per-AP regularized precoder is
\begin{align}
\label{eq:local-rzf}
\mat{W}_{l}[q] &= \big(\mat{E}_{l}[q]\mat{E}_{l}[q]^{H}+\lambda\mat{I}_{M}\big)^{-1}
\mat{E}_{l}[q] \inC{M\times K},\\
 l&=1,\dots,L,
\end{align}
whose column $k$ is the local direction $\bar{\vect{w}}_{kl}[q]$, and the
collective direction $\bar{\vect{w}}_{k}[q]$ stacks the $L$ per-AP blocks. Local
zero-forcing is the $\lambda=0$ limit, which nulls the interference an AP causes
to its own users only when $M\geq K$; because each AP is small (typically
$M<K$), the loading is needed to keep the $M\times M$ inverse well conditioned,
giving local \gls{rzf} for a heuristic $\lambda$ and local \gls{mmse} for
$\lambda=\sigma^{2}$. As in the centralized case, under perfect \gls{csi} local
\gls{mmse} coincides with local \gls{rzf} at that loading. The $M\times M$
antenna-domain form is used because each AP's array is small.

Because each AP designs and normalizes its own block independently, the power is
parametrized per user per AP: AP $l$ normalizes its power to respect $\rho_{kl}\geq 0$ for user $k$ this is achieved as
\begin{align}
    \vect{w}_{kl}[q]=\sqrt{\rho_{kl}}\,\bar{\vect{w}}_{kl}[q]/\sqrt{\expt{\norm{\bar{\vect{w}}_{kl}[q]}^{2}}}.
\end{align}
Fractional power control is used, each AP splits its own budget
$\rho_{\max}$ among its users according to their local large-scale gains,
\begin{equation}
\label{eq:fractional-power}
\rho_{kl} = \rho_{\max}\,
\frac{\beta_{kl}^{\,v}}{\sum_{i=1}^{K}\beta_{il}^{\,v}},
\qquad l=1,\dots,L,
\end{equation}
so that $\sum_{k=1}^{K}\rho_{kl}=\rho_{\max}$ by construction and every AP meets
its per-array budget in \eqref{eq:power-constraint} with equality, using only
local quantities and without any CPU-side global rescaling. The exponent
$v$ plays the same throughput-versus-fairness role as above, with
$v=1/2$ recovering the heuristic of \cite{demir2021foundations}. Beyond
such statistics-based rules the powers could be optimized for a network utility
such as max-min fairness or sum rate, but those problems do not scale to very
large networks, so the heuristics above are the practical default.
